\documentclass[10pt,a4paper]{article}
\usepackage[a4paper,top=1.5cm,bottom=1.4cm,left=1.7cm,right=1.7cm]{geometry}
\usepackage{fontspec}
\setmainfont{texgyreheros-regular.otf}[BoldFont=texgyreheros-bold.otf,ItalicFont=texgyreheros-italic.otf,BoldItalicFont=texgyreheros-bolditalic.otf]
\setmonofont{texgyrecursor-regular.otf}[BoldFont=texgyrecursor-bold.otf,Scale=MatchLowercase]
\usepackage{xcolor}
\usepackage{enumitem}
\usepackage{booktabs,array,tabularx}
\usepackage{tcolorbox}
\tcbuselibrary{skins}
\usepackage[hidelinks]{hyperref}
\definecolor{navy}{HTML}{131A2B}
\definecolor{accent}{HTML}{D95926}
\definecolor{blue}{HTML}{2563EB}
\definecolor{teal}{HTML}{0F766E}
\definecolor{soft}{HTML}{FFF4EC}
\pagestyle{empty}
\let\origund\_\renewcommand{\_}{\origund\allowbreak}
\sloppy
\setlength{\parindent}{0pt}
\setlength{\parskip}{3pt}
\usepackage{ragged2e}
\setlist{itemsep=2pt,topsep=2pt,parsep=0pt,leftmargin=12pt}
\newcommand{\hd}[1]{{\color{navy}\bfseries\large #1}\par\vspace{1pt}}
\newcommand{\lab}[1]{{\color{accent}\bfseries #1}}
\newtcolorbox{say}{colback=soft,colframe=accent,boxrule=0.6pt,arc=2pt,left=5pt,right=5pt,top=3pt,bottom=3pt,fontupper=\normalsize\setlength{\parskip}{4pt}}
\newtcolorbox{qa}{colback=white,colframe=navy!40,boxrule=0.5pt,arc=2pt,left=5pt,right=5pt,top=2pt,bottom=2pt,fontupper=\small}
\newcommand{\speaker}[3]{%
  \begin{tcolorbox}[colback=navy,colframe=navy,arc=3pt,left=8pt,right=8pt,top=4pt,bottom=4pt]
  {\color{white}\bfseries\Large #1}\hfill{\color{white!70!navy}\small #2}\\[1pt]{\color{accent!60!white}\bfseries #3}
  \end{tcolorbox}\vspace{-2pt}}

\begin{document}

% =============================================================== COVER
{\color{accent}\bfseries\Huge ReviewIQ}\\[2pt]
{\color{navy}\bfseries\Large Technical Viva Guide}\\[1pt]
{\small Vidyalankar Institute of Technology \quad|\quad Guide: Prof. Sneha Annappanavar}

\vspace{4pt}\hrule\vspace{6pt}

\hd{The 30-second pitch (anyone can say this)}
\begin{say}
Product reviews are unstructured text, and a star average hides detail. ReviewIQ takes 50,000 Amazon Electronics reviews and builds three layers on the same data: \textbf{NLP} to understand the language (sentiment, topics, aspects), \textbf{deep learning} to classify sentiment with a BiLSTM and a fine-tuned DistilBERT, and \textbf{advanced ML} to score helpfulness, cluster reviews and flag suspicious ones. A semantic-search question answering layer answers only from a product's own reviews and refuses when evidence is weak. Everything is in one Streamlit dashboard.
\end{say}

\hd{Who presents what (swap names freely)}
\begin{tabularx}{\textwidth}{@{}>{\bfseries\RaggedRight}p{4.6cm}>{\RaggedRight}p{3.6cm}>{\RaggedRight}X@{}}
\toprule
Speaker & Part & Covers \\
\midrule
Nidhi Bhawari (23102A0006) & 1. Data and NLP & Problem, dataset, labels, split, VADER, TF-IDF, LDA, aspect sentiment \\
Aarti Sakpal (23102A0007) & 2. Deep learning & BiLSTM with attention, DistilBERT, results, ablation, cost \\
Gayatri Patil (23102A0016) & 3. Advanced ML & XGBoost, KMeans, Isolation Forest \\
Prachi Wakode (23102A0020) & 4. Search, app, demo & Embeddings, RAG guardrail, dashboard, demo, limits \\
\bottomrule
\end{tabularx}

\hd{Order of the live demo (Prachi drives it, each speaker points at their part)}
\begin{enumerate}[leftmargin=16pt]
\item \textbf{Home:} three pillars and pipeline (everyone).
\item \textbf{Explore Products:} search ``echo'', open \emph{Echo Dot} with View AI Insights.
\item \textbf{Product Insights:} Aspect Analysis tab, click an aspect to show evidence sentences (Nidhi). Review Trust tab, drag the slider (Gayatri).
\item \textbf{Ask ReviewIQ:} ``Is this good for cutting cable TV?'' on the first Fire TV Stick, then ``Does it have a 3 year warranty?'' to show the refusal (Prachi).
\item \textbf{Deep Learning Lab, Try it live:} ``Mixed (aspect-level)'' example, click Classify review. The three models disagree and the attention words light up (Aarti).
\item \textbf{Business Intelligence, Complaint tab:} battery is the top complaint (Nidhi).
\end{enumerate}

\hd{Five numbers every speaker must know}
\begin{tabularx}{\textwidth}{@{}*{5}{>{\RaggedRight\small}X}@{}}
\textbf{\color{accent}50,000} reviews, 34,365 products & \textbf{\color{accent}88.8\%} / \textbf{0.718} BiLSTM accuracy / macro-F1 & \textbf{\color{accent}89.9\%} DistilBERT accuracy (6,000 train) & \textbf{\color{accent}57$\times$} LSTM faster on CPU, 26$\times$ fewer params & \textbf{\color{accent}20.2\%} battery mentions negative (top complaint) \\
\end{tabularx}

\hd{Golden rules for the viva}
\begin{itemize}
\item Say \textbf{``agreement with star ratings''}, not ``true accuracy'': the labels are weak (from stars).
\item Say \textbf{``potentially suspicious''}, never ``fake''. That module is unsupervised, so it has no accuracy figure.
\item If asked ``is the LSTM better than DistilBERT?'': better macro-F1 (0.718 vs 0.702) but \textbf{not statistically significant}; DistilBERT is better on accuracy (significant). DistilBERT saw only 6,000 reviews.
\item If you do not know a detail, say what the code does, not what you guess. Everything here is in \texttt{src/}.
\end{itemize}
\newpage

% =============================================================== SPEAKER 1
\speaker{Nidhi Bhawari}{23102A0006}{Part 1: Problem, data and NLP \quad (about 2.5 min)}
\begin{say}
\lab{Problem.} Reviews are free text, so a 4.4 star average cannot tell a buyer that the sound is great and the battery is poor. We ask three questions: should I buy this, what should be improved, can I trust these reviews.\par
\lab{Data.} Amazon Reviews 2023, Electronics: a 50,000 review sample, 34,365 products, 8,883 reviewers. Labels come from stars: 4 and 5 positive, 3 neutral, 1 and 2 negative, giving 80.4\%, 7.0\%, 12.6\%. Because it is so imbalanced we use \textbf{macro-F1}. I cleaned it (HTML and URLs stripped, fewer than 3 words dropped, duplicate user-product-text removed) and joined product name, brand and image by \texttt{parent\_asin} from a 5.2 GB metadata file (34,364 of 34,365 matched). Every model uses one \textbf{stratified 80/20 split, seed 42}: 40,000 train, 10,000 test.\par
\lab{NLP.} Four components. \textbf{VADER} is a lexicon scorer; compound $\ge0.05$ positive, $\le-0.05$ negative. \textbf{TF-IDF} with unigrams and bigrams, 30,000 features, sublinear tf ($1+\log \mathrm{tf}$), feeds a logistic regression and a linear SVM, both class-balanced. \textbf{LDA} finds 8 topics. \textbf{Aspect sentiment} finds what customers feel about each feature.
\end{say}
\par\vspace{2pt}\textbf{Technical detail to keep ready}\par
\begin{itemize}
\item \textbf{TF-IDF + SVM:} TF-IDF with 30,000 features, n-grams (1,2), \texttt{min\_df}=2, \texttt{sublinear\_tf}=True; \texttt{LinearSVC} with \texttt{class\_weight}=balanced; LR with \texttt{max\_iter}=1000, balanced.
\item \textbf{LDA:} count vectors with 3,000 features, bigrams, \texttt{min\_df}=5, \texttt{max\_df}=0.85, English stop words; online variational Bayes, 8 topics, 6 iterations, batch 1024; fitted on an 8,000-review subsample then applied to all 50,000 for speed.
\item \textbf{Aspect sentiment (ABSA):} sentence split with a regex on sentence-ending punctuation; 10 keyword lexicons (battery, camera, display, sound, performance, build quality, price, connectivity, size, customer service); the polarity of an aspect is VADER on the concatenation of the sentences that mention it; the matched sentence is stored as \textbf{evidence}. Result: 64,134 mentions in 33,257 reviews; \textbf{battery 20.2\% negative}, connectivity and customer service 19.3\%, build quality 8.1\%.
\item \textbf{Test results (10,000 reviews):} VADER .813 / .512, LR .860 / .689, SVM .888 / .693 (accuracy / macro-F1). VADER neutral F1 is only 0.08.
\end{itemize}
\par\vspace{2pt}\textbf{Likely questions}\par\vspace{1pt}
\begin{qa}\textbf{Why macro-F1 and not accuracy?} Predicting ``positive'' for everything already gives about 80\% accuracy. Macro-F1 averages the three class F1 scores equally, so the 7\% neutral class counts as much as the 80\% positive class.\end{qa}\vspace{2pt}
\begin{qa}\textbf{What is TF-IDF?} A word's weight is its frequency in the review times how rare it is across reviews, so common words matter less. Sublinear tf dampens repeated words.\end{qa}\vspace{2pt}
\begin{qa}\textbf{Why is the sample biased?} We took a contiguous block of the source file, which is grouped by reviewer (9,686 reviewers in 65,552 raw records), so prolific reviewers are over-represented. It matters for the suspicious-review feature.\end{qa}\vspace{2pt}
\begin{qa}\textbf{Is the aspect module validated?} Not against human labels. It is keyword matching plus VADER, so it misses implicit aspects and sarcasm. We say this in the limitations.\end{qa}\vspace{2pt}
\begin{qa}\textbf{Why does stratifying matter?} It keeps the 80/7/13 class ratio identical in train and test, so a rare class cannot vanish from the test set.\end{qa}
\newpage

% =============================================================== SPEAKER 2
\speaker{Aarti Sakpal}{23102A0007}{Part 2: Deep learning \quad (about 2.5 min)}
\begin{say}
\lab{BiLSTM with attention.} Text is lower-cased and tokenised with the regex \texttt{[a-z0-9']+ | [!?]}. The vocabulary is built on the training split only (14,937 tokens, minimum frequency 3) and reviews are truncated to 200 tokens. A learned 128-d embedding feeds a \textbf{2-layer bidirectional LSTM}, 128 units per direction, so each token gets a 256-d vector that sees both left and right context. \textbf{Attention pooling} scores every token, $\alpha_t=\mathrm{softmax}(\mathbf w^\top \mathbf h_t)$ with padding masked, and the sentence vector is $\mathbf c=\sum_t\alpha_t\mathbf h_t$. Then dropout 0.4 and a linear layer to 3 classes: \textbf{2,572,420 parameters}, about 7 minutes on an RTX 4050.\par
\lab{Imbalance.} The loss is cross-entropy with class weights $w_c=\sqrt{N/(3n_c)}$, and we stop on \textbf{validation macro-F1}, patience 3 (best epoch 6 of 9).\par
\lab{DistilBERT.} Pre-trained 6-layer transformer, 66.96 M parameters, fine-tuned with a 3-way head: AdamW $2\times10^{-5}$, 2 epochs, batch 8, max length 128, on a 6,000-review subsample on CPU (about 60 minutes).
\end{say}
\par\vspace{2pt}\textbf{Technical detail to keep ready}\par
\begin{itemize}
\item \textbf{Training:} Adam, lr $2\times10^{-3}$, weight decay $10^{-5}$, batch 128, gradient clipping 1.0, packed sequences (\texttt{pack\_padded\_\allowbreak sequence}) so padding is ignored by the LSTM. Split 36,000 train / 4,000 validation / 10,000 test.
\item \textbf{Results:} BiLSTM .888 / \textbf{.718}; DistilBERT \textbf{.899} / .702. Neutral F1: LSTM .446, SVM .382, DistilBERT .368.
\item \textbf{Significance (2,000 bootstrap resamples):} LSTM minus SVM macro-F1 $+0.025$, CI $[+0.010,+0.039]$, significant. LSTM minus DistilBERT $+0.016$, CI $[-0.000,+0.031]$, not significant. DistilBERT accuracy lead $+0.011$ is significant.
\item \textbf{Ablation (3 seeds):} full .891 / .722; no class weights .899 / .709; mean pooling .889 / .715. So class weights give $+0.012$ macro-F1 for $-0.007$ accuracy; attention adds only about $+0.007$. Seed spread is about 0.004.
\item \textbf{Cost:} 26$\times$ fewer parameters, 57$\times$ faster on CPU (975 vs 17 reviews/s), 13$\times$ on GPU.
\item \textbf{Word clouds:} word weight $=f_{class}\cdot(f_{class}/f_{all})^{0.7}$, so distinctive words win over merely frequent ones.
\end{itemize}
\par\vspace{2pt}\textbf{Likely questions}\par\vspace{1pt}
\begin{qa}\textbf{Why an LSTM at all?} It models word order and negation (``not good'') that bag-of-words cannot, is trained from scratch in minutes, and is cheap to serve. It is the course's required DL model.\end{qa}\vspace{2pt}
\begin{qa}\textbf{What does attention do?} It learns a weight per word so the sentence vector focuses on informative words. You can see it in the live demo: ``terrible'' and ``amazing'' are highlighted.\end{qa}\vspace{2pt}
\begin{qa}\textbf{Why $\sqrt{\cdot}$ in the class weight?} Full inverse-frequency weights over-correct and hurt accuracy; the square root is a milder compromise. We did not tune it, which we state.\end{qa}\vspace{2pt}
\begin{qa}\textbf{Is DistilBERT unfairly treated?} Partly: 6,000 reviews versus 36,000 for the LSTM, and no class weights. That is a stated limitation and the first item in future work.\end{qa}
\newpage

% =============================================================== SPEAKER 3
\speaker{Gayatri Patil}{23102A0016}{Part 3: Advanced machine learning \quad (about 2.5 min)}
\begin{say}
\lab{XGBoost helpfulness.} The target is whether a review got at least one helpful vote (26.2\% of reviews). Ten engineered features: rating, review length, title length, verified purchase, exclamation and question counts, capital-letter ratio, VADER compound, review age, and \textbf{number of aspects mentioned}, which comes from the NLP module. I tuned it with \texttt{RandomizedSearchCV}, 15 candidates, 3-fold, scored on F1, with \texttt{scale\_pos\_weight} for the imbalance. Best: 200 trees, depth 6, learning rate 0.01, subsample 0.7. Result: \textbf{70.5\% accuracy}, helpful-class precision .458, recall .702, F1 .555.\par
\lab{KMeans.} MiniBatchKMeans, $k=6$, on a separate TF-IDF space. Honest finding: only 2 clusters are topical, the other 4 are generic praise.\par
\lab{Isolation Forest.} Unsupervised, because we have no fake-review labels. It builds random trees; anomalies are isolated in few splits, giving score $s(x)=2^{-E[h(x)]/c(n)}$. Six features go through a StandardScaler; contamination 0.10 flags \textbf{4,999 reviews}.
\end{say}
\par\vspace{2pt}\textbf{Technical detail to keep ready}\par
\begin{itemize}
\item \textbf{XGBoost confusion matrix} (rows true, cols predicted): unhelpful 5211 / 2172, helpful 779 / 1838. Top features: review length, aspects mentioned, rating. Accuracy is modest because helpfulness depends on how many people saw the review, which no review feature captures.
\item \textbf{KMeans:} \texttt{MiniBatchKMeans(n\_clusters=6, n\_init=10, batch\_size=1024)} on a 5,000-feature TF-IDF with stop words removed. Sizes 19,309 / 3,570 / 13,552 / 3,228 / 5,101 / 5,240. Topical: laptop bags and cases, sound and headphones.
\item \textbf{Isolation Forest features:} is\_short ($<10$ words), capital ratio, exclamation count, verified flag, $\log(1+\mathrm{reviews})$, and max cosine similarity to another review of the \emph{same product}, computed per product group on the MiniLM embeddings (a full $50{,}000^2$ matrix would be 2.5 billion pairs). Score is $-$\texttt{score\_samples} min-max scaled to $[0,1]$; reasons are rule-based strings. 36 reviews have similarity $\ge0.95$.
\item \textbf{The bug we fixed:} the first version used raw rating and raw length; scores correlated $+0.55$ with length and $-0.42$ with rating, so long detailed negative reviews were flagged. After switching to directional features the correlations became $+0.06$ (length) and $+0.25$ (account volume).
\item \textbf{UI:} a threshold slider re-flags reviews live; the table shows flag, score, rating, review text and signals.
\end{itemize}
\par\vspace{2pt}\textbf{Likely questions}\par\vspace{1pt}
\begin{qa}\textbf{How do you know the flags are right?} We do not, and we say so. With no ground truth there is no accuracy; we inspected top flags (templated one-liners, ALL CAPS, high-volume unverified accounts) and call them ``potentially suspicious''.\end{qa}\vspace{2pt}
\begin{qa}\textbf{Why contamination 0.10?} It is an assumption about the anomaly rate, not a measurement; the slider lets users pick their own cut-off.\end{qa}\vspace{2pt}
\begin{qa}\textbf{Caveat on account volume?} The source file is reviewer-grouped, so counts per reviewer in our sample are inflated.\end{qa}\vspace{2pt}
\begin{qa}\textbf{What is \texttt{scale\_pos\_weight}?} Ratio of negative to positive training examples; it makes XGBoost penalise missing a helpful review more.\end{qa}
\newpage

% =============================================================== SPEAKER 4
\speaker{Prachi Wakode}{23102A0020}{Part 4: Search, question answering, dashboard and demo \quad (about 2.5 min)}
\begin{say}
\lab{Semantic search.} Every review is encoded once with \textbf{all-MiniLM-L6-v2} into a 384-d unit vector and stored, so cosine similarity is just a dot product. For a question about a product we mask to that product's reviews (\texttt{parent\_asin}) and take the \textbf{top 5}. \textbf{FLAN-T5-base} is prompted with the question and those reviews, truncated to 300 characters each, and answers in 1 to 2 sentences, greedy decoding.\par
\lab{The guardrail.} Small generators do not obey ``say you do not know'', so we added a hard rule: if the best similarity is below \textbf{0.30} the generator is never called and the answer is ``not enough information''. Measured over 9 products and 45 questions per group: \textbf{97.8\%} of answerable questions accepted, \textbf{95.6\%} of plausible but unanswerable and \textbf{100\%} of off-topic questions declined. Dense retrieval also beat TF-IDF on paraphrased aspect queries, 0.54 vs 0.32.\par
\lab{Dashboard.} Streamlit, nine pages. Data logic is in pure functions; heavy work is cached. Summaries and comparisons are \textbf{templates built from computed numbers, never an LLM}, so they cannot invent facts.
\end{say}
\par\vspace{2pt}\textbf{Technical detail to keep ready}\par
\begin{itemize}
\item \textbf{Why scope by product?} Otherwise reviews of product B can answer a question about product A. Retrieval is restricted to the selected product's reviews, so evidence from other products cannot appear.
\item \textbf{Caching:} \texttt{@st.cache\_data} for tables and features, \texttt{@st.cache\_resource} for models. Without it, the helpfulness view recomputed VADER on 50,000 reviews (about 23 s) on every click.
\item \textbf{Testing:} Streamlit \texttt{AppTest} runs all 9 pages for exceptions; Playwright clicked through 12 sidebar transitions and both in-page buttons.
\item \textbf{Bugs found and fixed:} backwards trust features, ungrounded answers, slow reruns, missing product images (wrong field name, dead fallback link), ignored first sidebar click (widget identity changed every page; fixed with a stable key and callbacks), and a folder named \texttt{pages} that Streamlit read as a second menu.
\item \textbf{Retrieval test caveat:} 10 queries, relevance judged by a keyword lexicon (so a lower bound); the battery query scored 0 for both methods.
\item \textbf{Limitations:} weak labels from stars; one English category in a reviewer-grouped sample; DistilBERT under-trained; keyword aspects; flags are anomalies; answer faithfulness was not measured.
\item \textbf{Future work:} DistilBERT on full data with matched class weights, learned aspect extraction, human-annotated subsets, random multi-category sampling, a faithfulness check for answers.
\end{itemize}
\par\vspace{2pt}\textbf{Likely questions}\par\vspace{1pt}
\begin{qa}\textbf{What is RAG?} Retrieve evidence first, then generate from it, so every answer can be traced to reviews. The user always sees the five retrieved reviews and their similarity scores.\end{qa}\vspace{2pt}
\begin{qa}\textbf{Why does the model refuse sometimes?} By design: when similarity is under 0.30 the evidence is too weak, which we measured separates answerable from unanswerable questions (mean 0.455 vs 0.218).\end{qa}\vspace{2pt}
\begin{qa}\textbf{Closing line:} ``Three families of AI, one corpus, every insight linked to evidence, and the limits stated honestly.''\end{qa}
\newpage

% =============================================================== CHEAT SHEET
{\color{navy}\bfseries\Large Technical cheat sheet}\hfill{\small keep this open during the viva}
\vspace{3pt}

\begin{tabularx}{\textwidth}{@{}>{\bfseries\RaggedRight}p{3.7cm}>{\RaggedRight}X@{}}
\toprule
Term & One-line meaning \\
\midrule
Macro-F1 & Mean of the per-class F1 scores, $F1_c=2P_cR_c/(P_c+R_c)$; weights rare and common classes equally \\
Stratified split & Same class ratio in train and test (80/7/13) \\
TF-IDF & Term frequency times inverse document frequency; down-weights common words \\
LDA & Probabilistic topic model: each review is a mixture of topics, each topic a distribution over words \\
VADER & Rule and lexicon sentiment scorer, compound score in $[-1,1]$ \\
LSTM / BiLSTM & Recurrent network with gates that keep long-range context; bidirectional reads both ways \\
Attention pooling & Learned weights over tokens: $\alpha_t=\mathrm{softmax}(\mathbf w^\top\mathbf h_t)$, $\mathbf c=\sum\alpha_t\mathbf h_t$ \\
Class weighting & Loss multiplied by $w_c=\sqrt{N/(3n_c)}$ so rare classes matter more \\
Early stopping & Stop when validation macro-F1 has not improved for 3 epochs, keep the best epoch \\
DistilBERT & Smaller transformer distilled from BERT; fine-tuned with a classification head \\
Bootstrap CI & Resample the test set 2,000 times; the 2.5th and 97.5th percentiles give the 95\% interval \\
XGBoost & Gradient-boosted decision trees; each tree fits the errors of the previous ones \\
Isolation Forest & Random trees isolate outliers in fewer splits; shorter path means more anomalous \\
Cosine similarity & $\cos(q,r)=\mathbf q\cdot\mathbf r$ for unit vectors; 1 is identical meaning \\
RAG & Retrieval-augmented generation: retrieve evidence, then generate from it \\
Similarity floor & If best $\cos<0.30$, refuse instead of generating \\
\bottomrule
\end{tabularx}

\vspace{4pt}
\par\textbf{All results on the same 10,000 test reviews}\par\vspace{2pt}
\begin{tabularx}{\textwidth}{@{}Xcccc@{}}
\toprule
Model & Accuracy & Macro-F1 & F1 neutral & Speed (CPU, reviews/s) \\
\midrule
VADER & .813 & .512 & .080 & n/a \\
TF-IDF + LR & .860 & .689 & .400 & n/a \\
TF-IDF + SVM & .888 & .693 & .382 & 13,550 \\
BiLSTM + Attention & .888 & \textbf{.718} & \textbf{.446} & 975 \\
DistilBERT & \textbf{.899} & .702 & .368 & 17 \\
\bottomrule
\end{tabularx}

\vspace{4pt}
\par\textbf{Files to know} (all under \texttt{src/})\par\vspace{2pt}
\begin{tabularx}{\textwidth}{@{}>{\ttfamily\footnotesize\RaggedRight}p{6.4cm}>{\RaggedRight}X@{}}
data\_prep.py, metadata\_prep.py & Cleaning, labels, product metadata join \\
sentiment/, topic\_modeling/, absa/ & VADER, TF-IDF models, LDA, aspect sentiment \\
deep\_learning/lstm\_sentiment.py, lstm\_model.py & BiLSTM model and training \\
deep\_learning/bert\_finetune.py & DistilBERT fine-tuning \\
advanced\_ml/ & XGBoost helpfulness, KMeans, Isolation Forest \\
semantic\_search/, rag\_chatbot/ & Embeddings, retrieval, guardrail, answers \\
dashboard/ (app.py, sections/) & The 9-page Streamlit app \\
evaluation/ & Bootstrap, ablation, cost, retrieval and guardrail experiments \\
\bottomrule
\end{tabularx}

\end{document}
